% Image credits for the photographs and AI illustrations of Book 3
% (University Chemistry, Year 2), Hindi edition. Machine-readable license
% records live in images/book3/CREDITS.md; keep this file in sync with the
% English frontmatter/image-credits-book3.tex. Authors' names stay in their
% original Latin spelling here (attribution), as in the Book 2 Hindi credits.
\chapter*{चित्र-श्रेय}

इस पुस्तक के रेखाचित्र मौलिक हैं। जहाँ छायाचित्र प्रयुक्त हैं, वे अपने-अपने मुक्त
अनुज्ञापत्रों के अंतर्गत प्रयुक्त हैं, उनके रचयिताओं के प्रति आभार सहित; हर छायाचित्र के
स्रोत की पूरी कड़ियाँ और अनुज्ञापत्रों के पते पुस्तक के रिपॉज़िटरी में
\texttt{images/book3/CREDITS.md} में सूचीबद्ध हैं।
संकट चित्रलेख रसायनों के वर्गीकरण और लेबलन की विश्व स्तर पर सुसंगत प्रणाली (GHS) के
चित्रलेख हैं, जैसे उन्हें \TeX\ Live का \texttt{ghsystem} पैकेज खींचता है।

\bigskip
छायाचित्रों और मौलिक रेखाचित्रों के साथ-साथ, कुछ चित्र
कृत्रिम बुद्धि से बनाए गए चित्रण हैं, जो पुस्तक के संपादकों द्वारा लिखे गए निर्देशों से
OpenAI के चित्र-निर्माण द्वारा बनाए गए हैं, और समावेश से पहले रासायनिक
यथार्थता के लिए जाँचे गए हैं। हर निर्मित चित्र का पूरा निर्देश
रिपॉज़िटरी में \texttt{images/book3/ai/PROMPTS.md} में अभिलिखित है।

\bigskip
\noindent छायाचित्र, अध्याय के अनुसार:
\begin{itemize}\small
  \item अध्याय~1: जर्मेन हेनरी हेस, I.~I. Reimers द्वारा चित्र (1837--1838),
        CC BY-SA 4.0; मार्सेलें बर्थेलो, अज्ञात लेखक, सार्वजनिक अधिकार-क्षेत्र
        (दोनों Wikimedia Commons)।
  \item अध्याय~2: जोसाया विलर्ड गिब्स, अज्ञात छायाकार, सार्वजनिक अधिकार-क्षेत्र
        (Wikimedia Commons)।
  \item अध्याय~3: फ़्रांस्वा-मारी राउल्ट, अज्ञात छायाकार, सार्वजनिक अधिकार-क्षेत्र
        (Wikimedia Commons)।
  \item अध्याय~4: याकोबस हेंड्रिकस वांट हॉफ़, Nicola Perscheid द्वारा (1904),
        सार्वजनिक अधिकार-क्षेत्र; हेनरी ले शातेलिए, अज्ञात लेखक, CC BY-SA 4.0;
        फ़्रिट्ज़ हेबर और कार्ल बॉश, Nobel Foundation, सार्वजनिक अधिकार-क्षेत्र (सभी
        Wikimedia Commons)।
  \item अध्याय~6: डुइसबुर्ग-नॉर्ड की वात्या भट्टी, Dietmar Rabich द्वारा, CC BY-SA
        4.0; रेल की थर्माइट वेल्डिंग, PetrS. द्वारा, CC BY-SA 3.0 (दोनों
        Wikimedia Commons)।
  \item अध्याय~7: एक रिफ़ाइनरी का आसवन स्तंभ, CEphoto, Uwe Aranas द्वारा,
        CC BY-SA 3.0 (Wikimedia Commons)।
  \item अध्याय~9: माइकल फ़ैराडे, Thomas Phillips द्वारा (1842), सार्वजनिक अधिकार-क्षेत्र
        (Wikimedia Commons)।
  \item अध्याय~10: यारोस्लाव हेयरोव्स्की, J.~Heyrovsk\'y Institute of Physical
        Chemistry का अभिलेखागार, CC BY-SA 3.0 (Wikimedia Commons)।
  \item अध्याय~11: एक वोल्टीय पुंज, GuidoB द्वारा, CC BY-SA 3.0; चार्ल्स मार्टिन हॉल
        और पॉल एरू, अज्ञात लेखक, सार्वजनिक अधिकार-क्षेत्र (सभी Wikimedia
        Commons)।
  \item अध्याय~12: एक ताम्र प्रतिमा, Elcobbola द्वारा, सार्वजनिक अधिकार-क्षेत्र; एक बलिदानी
        ऐनोड, Zwergelstern द्वारा, CC BY-SA 3.0 (दोनों Wikimedia Commons)।
  \item अध्याय~13: एर्विन श्रोडिंगर, Nobel Foundation (1933), सार्वजनिक अधिकार-क्षेत्र
        (Wikimedia Commons)।
  \item अध्याय~14: चुंबक में द्रव ऑक्सीजन, Pieter Kuiper द्वारा, सार्वजनिक अधिकार-क्षेत्र;
        रॉबर्ट मुलिकेन और फ़्रीडरिख हुंड, GFHund द्वारा, CC BY 3.0 (दोनों Wikimedia
        Commons)।
  \item अध्याय~16: आउगुस्ट केकुले (1873), अज्ञात लेखक, सार्वजनिक अधिकार-क्षेत्र
        (Wikimedia Commons)।
  \item अध्याय~18: संक्रमण धातु विलयन, Benjah-bmm27 द्वारा, सार्वजनिक अधिकार-क्षेत्र; आल्फ़्रेड
        वर्नर, Les Prix Nobel (1914), सार्वजनिक अधिकार-क्षेत्र (दोनों Wikimedia Commons)।
  \item अध्याय~19: माणिक का एक क्रिस्टल, सार्वजनिक अधिकार-क्षेत्र; एक पन्ना, Didier Descouens द्वारा,
        CC BY-SA 3.0; कॉपर सल्फ़ेट के क्रिस्टल, Stephanb द्वारा, CC BY-SA 3.0 (सभी
        Wikimedia Commons)।
  \item अध्याय~20: आकिरा सुज़ुकी, एइइचि नेगिशी और रिचर्ड हेक, BloodIce द्वारा, CC BY-SA
        4.0 (Wikimedia Commons)।
  \item अध्याय~22: शार्ल फ़्रीडेल (1890 का दशक) और जेम्स मेसन क्राफ़्ट्स (\textit{Popular
        Science Monthly}, 1900), सार्वजनिक अधिकार-क्षेत्र (Wikimedia Commons)।
  \item अध्याय~24: मिशेल यूजीन शेव्रोल, Nicolas Eustache Maurin द्वारा, सार्वजनिक अधिकार-क्षेत्र
        (Wikimedia Commons)।
  \item अध्याय~25: अलेक्सांद्र बोरोदिन, अज्ञात लेखक, सार्वजनिक अधिकार-क्षेत्र; लुडविग क्लाइज़न,
        Veronica.villagrasa द्वारा, CC BY-SA 4.0 (दोनों Wikimedia Commons)।
  \item अध्याय~26: रॉबर्ट रॉबिन्सन, अज्ञात लेखक, सार्वजनिक अधिकार-क्षेत्र (Wikimedia Commons)।
  \item अध्याय~28: इलायस जेम्स कोरी, Trvthchem द्वारा, CC0 (Wikimedia Commons)।
  \item अध्याय~29: प्रयोगशाला में वॉलेस कैरोदर्स, अज्ञात छायाकार, सार्वजनिक अधिकार-क्षेत्र
        (Wikimedia Commons)।
  \item अध्याय~30: एमिल फ़िशर, Atelier Victoria, बर्लिन, लगभग 1895, सार्वजनिक अधिकार-क्षेत्र (Wikimedia
        Commons)।
  \item अध्याय~31: मिख़ाइल त्स्वेत, अज्ञात लेखक, सार्वजनिक अधिकार-क्षेत्र (Wikimedia Commons)।
  \item अध्याय~32: फ़्रांसिस विलियम ऐस्टन (1922), अज्ञात लेखक, सार्वजनिक अधिकार-क्षेत्र; गुस्ताफ़ किरखॉफ़,
        रॉबर्ट बुन्सन और हेनरी रॉस्को (1862), सार्वजनिक अधिकार-क्षेत्र (दोनों Wikimedia Commons)।
  \item अध्याय~34: विलियम सीली गॉसेट (1908), सार्वजनिक अधिकार-क्षेत्र (Wikimedia Commons)।
%% next chapter here
\end{itemize}
\clearpage
